\begin{corollary}[pi-alt]
\label{pi_alt}
\lean{pi_alt}
% at /home/runner/work/trough/trough/library/PrimeNumberTheoremAnd/Consequences.lean:853.0-898.9
\leanok
One has
  $$ \pi(x) = (1+o(1)) \frac{x}{\log x}$$
  as $x \to \infty$.
\end{corollary}
\begin{proof}
\uses{pi_asymp}
\leanok
An integration by parts gives
  $$ \int_2^x \frac{dt}{\log t} = \frac{x}{\log x} - \frac{2}{\log 2} +
  \int_2^x \frac{dt}{\log^2 t}.$$
  We have the crude bounds
  $$ \int_2^{\sqrt{x}} \frac{dt}{\log^2 t} = O( \sqrt{x} )$$
  and
  $$ \int_{\sqrt{x}}^x \frac{dt}{\log^2 t} = O( \frac{x}{\log^2 x} )$$
  and combining all this we obtain
  $$ \int_2^x \frac{dt}{\log t} = \frac{x}{\log x} + O( \frac{x}{\log^2 x} )$$
  $$ = (1+o(1)) \frac{x}{\log x}$$
  and the claim then follows from Theorem \ref{pi_asymp}.
\end{proof}
